% GENERATED FILE. Edit canonical lessons, then run npm run generate:books.
% canonical-source-hash: fnv1a64:0baa3eef88fbbd98
% canonical-lessons: MW-C105-raivno, MW-C105-bachno, MW-C105-janno, MW-C105-bharno, MW-C105-khichno

\chapter{To stay, To read, To know, To fill, To pull}
\label{ch:mw-to-stay-to-read-to-know-to-fill-to-pull}
\hlchaptermodality{105}

\begin{chapteropening}
\textbf{By the end of this chapter:} \emph{I can say to stay, to read, to know, to fill and to pull.}

Complete the last lesson of chapter 105: i can say to stay, to read, to know, to fill and to pull.
\end{chapteropening}

\section[\texorpdfstring{\mw{रैवणो} (raivṇo)}{रैवणो (raivṇo)}]{\mw{रैवणो} (raivṇo) — to stay, to live}

\label{lesson:MW-C105-raivno}



\begin{quote}
Before the new one: say the Marwadi for to ask for, then the Marwadi for to look.
\end{quote}

\subsection*{You'll want to know: \mw{रैवणो}}
\textbf{\mw{रैवणो}} — \emph{raivṇo} — \textquotedblleft{}to stay, to live\textquotedblright{}.

\begin{grammarlens}[title={What you've built}]
One more word for this chapter.
\end{grammarlens}

\subsection*{Your turn}
\begin{itemize}
\raggedright
  \item \emph{Say it:} \emph{raivṇo}
  \item \emph{Say it:} \emph{raivṇo}, once more
  \item \emph{Say it:} say \emph{raivṇo}
  \item \emph{Recall:} say the Marwadi for to ask for, then the Marwadi for to look, then say \emph{raivṇo} again
\end{itemize}

\begin{tcolorbox}[breakable,colback=teal!4,colframe=teal!35!black,title={Before you move on}]
What does \mw{रैवणो} mean? (\textquotedblleft{}To stay, to live\textquotedblright{}.) Say it once more.
\end{tcolorbox}
\section[\texorpdfstring{\mw{बांचणो} (bā̃chṇo)}{बांचणो (bā̃chṇo)}]{\mw{बांचणो} (bā̃chṇo) — to read}

\label{lesson:MW-C105-bachno}



\begin{quote}
Before the new one: say the Marwadi for to look, then the Marwadi for to stay.
\end{quote}

\subsection*{You'll want to know: \mw{बांचणो}}
\textbf{\mw{बांचणो}} — \emph{bā̃chṇo} — \textquotedblleft{}to read\textquotedblright{}.

\begin{grammarlens}[title={What you've built}]
One more word for this chapter.
\end{grammarlens}

\subsection*{Your turn}
\begin{itemize}
\raggedright
  \item \emph{Say it:} \emph{bā̃chṇo}
  \item \emph{Say it:} \emph{bā̃chṇo}, once more
  \item \emph{Say it:} say \emph{bā̃chṇo}
  \item \emph{Recall:} say the Marwadi for to look, then the Marwadi for to stay, then say \emph{bā̃chṇo} again
\end{itemize}

\begin{tcolorbox}[breakable,colback=teal!4,colframe=teal!35!black,title={Before you move on}]
What does \mw{बांचणो} mean? (\textquotedblleft{}To read\textquotedblright{}.) Say it once more.
\end{tcolorbox}
\section[\texorpdfstring{\mw{जाणणो} (jāṇṇo)}{जाणणो (jāṇṇo)}]{\mw{जाणणो} (jāṇṇo) — to know}

\label{lesson:MW-C105-janno}



\begin{quote}
Before the new one: say the Marwadi for to stay, then the Marwadi for to read.
\end{quote}

\subsection*{You'll want to know: \mw{जाणणो}}
\textbf{\mw{जाणणो}} — \emph{jāṇṇo} — \textquotedblleft{}to know\textquotedblright{}.

\begin{grammarlens}[title={What you've built}]
One more word for this chapter.
\end{grammarlens}

\subsection*{Your turn}
\begin{itemize}
\raggedright
  \item \emph{Say it:} \emph{jāṇṇo}
  \item \emph{Say it:} \emph{jāṇṇo}, once more
  \item \emph{Say it:} say \emph{jāṇṇo}
  \item \emph{Recall:} say the Marwadi for to stay, then the Marwadi for to read, then say \emph{jāṇṇo} again
\end{itemize}

\begin{tcolorbox}[breakable,colback=teal!4,colframe=teal!35!black,title={Before you move on}]
What does \mw{जाणणो} mean? (\textquotedblleft{}To know\textquotedblright{}.) Say it once more.
\end{tcolorbox}
\section[\texorpdfstring{\mw{भरणो} (bharṇo)}{भरणो (bharṇo)}]{\mw{भरणो} (bharṇo) — to fill}

\label{lesson:MW-C105-bharno}



\begin{quote}
Before the new one: say the Marwadi for to read, then the Marwadi for to know.
\end{quote}

\subsection*{You'll want to know: \mw{भरणो}}
\textbf{\mw{भरणो}} — \emph{bharṇo} — \textquotedblleft{}to fill\textquotedblright{}.

\begin{grammarlens}[title={What you've built}]
One more word for this chapter.
\end{grammarlens}

\subsection*{Your turn}
\begin{itemize}
\raggedright
  \item \emph{Say it:} \emph{bharṇo}
  \item \emph{Say it:} \emph{bharṇo}, once more
  \item \emph{Say it:} say \emph{bharṇo}
  \item \emph{Recall:} say the Marwadi for to read, then the Marwadi for to know, then say \emph{bharṇo} again
\end{itemize}

\begin{tcolorbox}[breakable,colback=teal!4,colframe=teal!35!black,title={Before you move on}]
What does \mw{भरणो} mean? (\textquotedblleft{}To fill\textquotedblright{}.) Say it once more.
\end{tcolorbox}
\section[\texorpdfstring{\mw{खींचणो} (khī̃chṇo)}{खींचणो (khī̃chṇo)}]{\mw{खींचणो} (khī̃chṇo) — to pull}

\label{lesson:MW-C105-khichno}



\begin{quote}
Before the new one: say the Marwadi for to know, then the Marwadi for to fill.
\end{quote}

\subsection*{You'll want to know: \mw{खींचणो}}
\textbf{\mw{खींचणो}} — \emph{khī̃chṇo} — \textquotedblleft{}to pull\textquotedblright{}.

\begin{grammarlens}[title={What you've built}]
One more word for this chapter.
\end{grammarlens}

\subsection*{Your turn}
\begin{itemize}
\raggedright
  \item \emph{Say it:} \emph{khī̃chṇo}
  \item \emph{Say it:} \emph{khī̃chṇo}, once more
  \item \emph{Say it:} say \emph{khī̃chṇo}
  \item \emph{Recall:} say the Marwadi for to know, then the Marwadi for to fill, then say \emph{khī̃chṇo} again
\end{itemize}

\begin{tcolorbox}[breakable,colback=teal!4,colframe=teal!35!black,title={Before you move on}]
What does \mw{खींचणो} mean? (\textquotedblleft{}To pull\textquotedblright{}.) Say it once more.
\end{tcolorbox}
